% ==================================================================
%  徐森林、薛春华《数学分析》(第一册)
%  第 6 章复习题
% ==================================================================
%  源: 原书印刷页 444-449(PDF 物理页 450-455)
%  底本: PaddleOCR-VL-1.6 + 高清影印逐题核验校订
%  规范: exblock 容器; 全西文半角标点; 单一落盘出口 _texout.py
% ==================================================================
\begin{exblock}{复习题 6}

\item (1) 设 $f$ 在 $[a,b]$ 上 2 阶可导, 且 $f''(x)>0$. 证明:
① $f\left(\frac{a+b}{2}\right)\leqslant\frac{1}{b-a}\int_{a}^{b}f(x)\mathrm{d}x$;
② 若 $f(x)\leqslant 0$, $x\in[a,b]$, 则有
\[
f(x)\geqslant\frac{2}{b-a}\int_{a}^{b}f(x)\mathrm{d}x,\quad x\in[a,b];
\]
(2) 设凸函数 $f$ 在 $[a,b]$ 上 Riemann 可积, 证明
\[
f\left(\frac{a+b}{2}\right)\leqslant\frac{1}{b-a}\int_{a}^{b}f(x)\mathrm{d}x\leqslant\frac{f(a)+f(b)}{2}.
\]
(3) 设 $\varphi$ 在 $[0,a]$ 上连续, $f(x)$ 2 阶可导, 且 $f''(x)\geqslant 0$, 证明:
\[
\frac{1}{a}\int_{0}^{a}f(\varphi(t))\mathrm{d}t\geqslant f\left(\frac{1}{a}\int_{0}^{a}\varphi(t)\mathrm{d}t\right).
\]
(4) 设 $f$ 在 $[a,b]$ 上连续, 且 $f(x)>0$, 则
\[
\ln\left(\frac{1}{b-a}\int_{a}^{b}f(x)\mathrm{d}x\right)\geqslant\frac{1}{b-a}\int_{a}^{b}\ln f(x)\mathrm{d}x.
\]

\item (1) 设 $f(x)$ 在 $[a,b]$ 上具有 1 阶导函数, $f(a)=f(b)=0$. 证明: $\exists\,\xi\in(a,b)$, s.t.
\[
|f'(\xi)|\geqslant\frac{4}{(b-a)^{2}}\left|\int_{a}^{b}f(x)\mathrm{d}x\right|.
\]
(2) 设 $f(x)$ 在 $[a,b]$ 上具有 1 阶导函数, $f(a)=f(b)=0$, 且 $f(x)\not\equiv 0$. 证明: $\exists\,\xi\in(a,b)$, s.t.
\[
|f'(\xi)|>\frac{4}{(b-a)^{2}}\int_{a}^{b}|f(x)|\mathrm{d}x\geqslant\frac{4}{(b-a)^{2}}\left|\int_{a}^{b}f(x)\mathrm{d}x\right|.
\]
(3) 设 $f(x)$ 在 $[a,b]$ 上具有 1 阶导函数, $f(a)=f(b)=0$, 且 $f(x)\not\equiv 0$, 则 $\exists\,\xi\in(a,b)$, s.t.
\[
|f'(\xi)|>\frac{4}{(b-a)^{2}}\left|\int_{a}^{b}f(x)\mathrm{d}x\right|
\]
等价于设 $F(x)$ 在 $[a,b]$ 上 2 阶可导, $F'(a)=F'(b)=0$, 且 $F(x)\not\equiv$ 常数, 则 $\exists\,\xi\in(a,b)$, s.t.
\[
|F''(\xi)|>\frac{4}{(b-a)^{2}}|F(b)-F(a)|.
\]

\item (1) 设 $f$ 在 $[a,b]$ 上可导, $f\left(\frac{a+b}{2}\right)=0$. 证明: $\exists\,\xi\in(a,b)$, s.t.
\[
|f'(\xi)|\geqslant\frac{4}{(b-a)^{2}}\int_{a}^{b}|f(x)|\mathrm{d}x\geqslant\frac{4}{(b-a)^{2}}\left|\int_{a}^{b}f(x)\mathrm{d}x\right|.
\]
(2) 设 $f$ 在 $[a,b]$ 上可导, $f\left(\frac{a+b}{2}\right)=0$, $f(x)\not\equiv 0$. 证明: $\exists\,\xi\in(a,b)$, s.t.
\[
|f'(\xi)|>\frac{4}{(b-a)^{2}}\left|\int_{a}^{b}f(x)\mathrm{d}x\right|.
\]
举例说明: $|f'(\xi)|>\frac{4}{(b-a)^{2}}\int_{a}^{b}|f(x)|\mathrm{d}x$ 未必成立.
(3) 上述 (2) $\Leftrightarrow$ 设 $F(x)$ 在 $[a,b]$ 上 2 阶可导, $F'\left(\frac{a+b}{2}\right)=0$, $F(x)\not\equiv$ 常数, 则 $\exists\,\xi\in(a,b)$, s.t.
\[
|F''(\xi)|>\frac{4}{(b-a)^{2}}|F(b)-F(a)|.
\]

\item 设 $a,b>0$, $f$ 在 $[-a,b]$ 上连续. 又设 $f>0$ 且 $\int_{-a}^{b}xf(x)\mathrm{d}x=0$. 证明:
\[
\int_{-a}^{b}x^{2}f(x)\mathrm{d}x\leqslant ab\int_{-a}^{b}f(x)\mathrm{d}x.
\]

\item 证明: 不等式
\[
\int_{-1}^{1}(1-x^{2})^{n}\mathrm{d}x\geqslant\frac{4}{3\sqrt{n}},\quad n\in\mathbb{N},
\]
且当 $n\geqslant 2$ 时严格不等号成立.

\item 证明: $\int_{0}^{\pi}\cos nx\cos^{n}x\,\mathrm{d}x=\frac{\pi}{2^{n}}$, $n=0,1,2,\cdots$.

\item 设 $f$ 在 $[0,1]$ 上连续, 并且 $\int_{0}^{1}x^{k}f(x)\,\mathrm{d}x=0$, $k=0,1,2,\cdots,n-1$, $\int_{0}^{1}x^{n}f(x)\,\mathrm{d}x=1$. 证明: $\exists\,\xi\in(0,1)$, s.t. $|f(\xi)|\geqslant 2^{n}(n+1)$.

\item 设函数 $f$ 在 $[a,b]$ 上连续、非负且严格单调增, 由积分中值定理, $\forall n\in\mathbb{N}$, $\exists_{1}\,x_{n}\in[a,b]$, s.t.
\[
f^{n}(x_{n})=\frac{1}{b-a}\int_{a}^{b}f^{n}(t)\mathrm{d}t.
\]
证明: $\lim\limits_{n\to+\infty}x_{n}=b$.

\item 设 $f$ 是一个 $n$ 次多项式, 且满足 $\int_{0}^{1}x^{k}f(x)\,\mathrm{d}x=0$, $k=1,2,\cdots,n$. 证明:
\[
\int_{0}^{1}f^{2}(x)\mathrm{d}x=(n+1)^{2}\left(\int_{0}^{1}f(x)\mathrm{d}x\right)^{2}.
\]

\item 设函数 $f$ 在区间 $[0,+\infty)$ 上单调增, 证明: $\varphi(x)=\int_{0}^{x}f(t)\mathrm{d}t$ 在 $[0,+\infty)$ 上为凸函数. 如果 $f$ 在区间 $[0,+\infty]$ 上单调增且连续, 试用简单方法证明之.

\item 设 $f(x)$ 2 阶连续可导, $f\geqslant 0$, $f'\leqslant 0$. 证明: $\forall c\in[a,b]$, 有
\[
f(c)\leqslant\frac{2}{b-a}\int_{a}^{b}f(x)\mathrm{d}x.
\]

\item 对 $n\in\mathbb{N}$, 定义 $I_{n}=\int_{0}^{\frac{\pi}{2}}\frac{\sin^{2}nt}{\sin t}\,\mathrm{d}t$, 求极限 $\lim\limits_{n\to+\infty}\frac{I_{n}}{\ln n}$.

\item 设 $f$ 在 $[a,b]$ 上连续, 若 $\int_{a}^{b}f(x)g(x)\,\mathrm{d}x=0$ 对 $[a,b]$ 上一切满足条件 $g(a)=g(b)=0$ 的连续函数 $g$ 成立. 证明: 在 $[a,b]$ 上, $f(x)\equiv 0$.

\item 设函数 $f$ 在 $[-1,1]$ 上可导, $M=\sup|f'|$. 如果 $\exists\,a\in(0,1)$, s.t. $\int_{-a}^{a}f(x)\,\mathrm{d}x=0$. 证明:
\[
\left|\int_{-1}^{1}f(x)\mathrm{d}x\right|\leqslant M(1-a^{2}).
\]

\item 设 $f$ 在 $[a,b]$ 上连续可导. 证明:
\[
\max_{a\leqslant x\leqslant b}|f(x)|\leqslant\frac{1}{b-a}\left|\int_{a}^{b}f(x)\mathrm{d}x\right|+\int_{a}^{b}|f'(x)|\,\mathrm{d}x.
\]

\item 设 $f$ 在 $(0,+\infty)$ 上连续, 且 $\forall x>0$ 有
\[
\int_{x}^{x^{2}}f(t)\,\mathrm{d}t=\int_{1}^{x}f(t)\,\mathrm{d}t,
\]
试求出满足上述条件的一切函数 $f$.

\item 设 $f$ 在 $[0,1]$ 上 2 阶连续可导, $f(0)=f(1)=0$, 且当 $x\in(0,1)$ 时, $f(x)\neq 0$. 证明:
\[
\int_{0}^{1}\left|\frac{f''(x)}{f(x)}\right|\mathrm{d}x\geqslant 4.
\]

\item 设 $f$ 在 $[a,b]$ 上连续、单调增. 证明:
\[
\int_{a}^{b}xf(x)\,\mathrm{d}x\geqslant\frac{a+b}{2}\int_{a}^{b}f(x)\,\mathrm{d}x.
\]
(提示: 令 $F(t)=\int_{a}^{t}xf(x)\,\mathrm{d}x-\frac{a+t}{2}\int_{a}^{t}f(x)\,\mathrm{d}x$.)

\item 设 $b>a>0$, 证明不等式:
\[
\ln\frac{b}{a}>\frac{2(b-a)}{a+b}.
\]
(提示: 对 $f(x)=-\frac{1}{x}$ 应用题 18 的结果.)

\item 设函数 $f$ 在 $[0,1]$ 上 2 阶连续可导, 且 $f(0)=f(1)=f'(0)=0$, $f'(1)=1$. 证明:
\[
\int_{0}^{1}\left[f''(x)\right]^{2}\mathrm{d}x\geqslant 4.
\]
指出式中等号成立的条件.

\item 设函数 $f$ 在 $[0,1]$ 上 Riemann 可积, 且有正数 $m$ 与 $M$, 使得 $m\leqslant f(x)\leqslant M$, $\forall x\in[0,1]$. 证明:
\[
1\leqslant\int_{0}^{1}f(x)\,\mathrm{d}x\int_{0}^{1}\frac{\mathrm{d}x}{f(x)}\leqslant\frac{(m+M)^{2}}{4mM}.
\]

\item 设 $x(t)$ 在 $[0,a]$ 上连续且满足
\[
|x(t)|\leqslant M+k\int_{0}^{t}|x(t)|\,\mathrm{d}t,
\]
这里 $M$ 与 $k$ 为正常数. 证明: $|x(t)|\leqslant M\mathrm{e}^{kt}$, $t\in[0,a]$.

\item 设 $f$ 为 $[-1,1]$ 上的连续函数, 且对 $[-1,1]$ 上的任何偶函数 $g$, 积分 $\int_{-1}^{1}f(x)g(x)\,\mathrm{d}x=0$. 证明: $f$ 在 $[-1,1]$ 上为奇函数.

\item 设 $a,b,n\in\mathbb{N}$, 令 $f(x)=\frac{x^{n}(a-bx)^{n}}{n!}$. 证明:
\begin{enumerate}[label=(\arabic*)]
  \item $f\left(\frac{a}{b}-x\right)=f(x)$;
  \item $f^{(k)}(x)$ ($0\leqslant k\leqslant 2n$) 当 $x=0$, $x=\frac{a}{b}$ 时取得值为整数;
  \item 假设 $\pi$ 为有理数, 即 $\pi=\frac{a}{b}$, $a$ 与 $b$ 为既约正整数. 则可证 $\int_{0}^{\pi}f(x)\sin x\,\mathrm{d}x$ 为整数. 由此证明 $\pi$ 不可能为有理数, 即 $\pi$ 为无理数.
\end{enumerate}

\item 设函数 $f$ 在 $[0,2]$ 上连续可导, $f(0)=f(2)=1$, 且 $|f'|\leqslant 1$. 证明: $1\leqslant\int_{0}^{2}f(x)\,\mathrm{d}x\leqslant 3$.

\item 证明: $\int_{0}^{x}\mathrm{e}^{xt-t^{2}}\mathrm{d}t=\mathrm{e}^{\frac{x^{2}}{4}}\int_{0}^{x}\mathrm{e}^{-\frac{t^{2}}{4}}\mathrm{d}t$.

\item 证明: (1) $\forall n\in\mathbb{N}$, 积分
\[
I_{n}=\int_{0}^{1}\frac{1}{x}\left[\left(\ln x+\ln\frac{1+x}{1-x}\right)^{n}-\ln^{n}x\right]\mathrm{d}x
\]
是收敛的;
(2) 当 $n$ 为偶数时, $I_{n}=0$.

\item 证明: $\int_{0}^{+\infty}\left(\frac{\sin x}{x}\right)^{2}\mathrm{d}x=\int_{0}^{+\infty}\frac{\sin x}{x}\mathrm{d}x$.

\item 计算积分:
\begin{enumerate}[label=(\arabic*)]
  \item $\int_{0}^{+\infty}\frac{\ln x}{(1+x)\sqrt{x}}\,\mathrm{d}x$;
  \item $\int_{0}^{\frac{\pi}{2}}\frac{\mathrm{d}x}{1+\tan^{100}x}$;
  \item $\int_{1}^{+\infty}\left(\frac{1}{[x]}-\frac{1}{x}\right)\mathrm{d}x$;
  \item $\int_{0}^{+\infty}\frac{\mathrm{e}^{-x^{2}}}{\left(x^{2}+\frac{1}{2}\right)^{2}}\,\mathrm{d}x$.
\end{enumerate}

\item 设 $f$ 为 $[-1,1]$ 上的连续函数. 证明:
\[
\lim_{h\to 0^{+}}\int_{-1}^{1}\frac{h}{h^{2}+x^{2}}f(x)\,\mathrm{d}x=\pi f(0).
\]

\item 设 $f$ 为 $[-1,1]$ 上的连续函数. 证明:
\[
\lim_{\lambda\to 0^{+}}\frac{1}{2\lambda}\int_{-1}^{1}f(x)\mathrm{e}^{-\frac{|x|}{\lambda}}\,\mathrm{d}x=f(0).
\]

\item 设 $f$ 在 $[0,\pi]$ 上 Riemann 可积. 证明: 不能同时有
\[
\int_{0}^{\pi}|f(x)-\sin x|^{2}\,\mathrm{d}x\leqslant\frac{3}{4},\quad\int_{0}^{\pi}|f(x)-\cos x|^{2}\,\mathrm{d}x\leqslant\frac{3}{4}.
\]

\item 设 $f$ 在 $[a,b]$ 上 2 阶连续可导. 证明: $\exists\,\xi\in(a,b)$, s.t.
\[
\int_{a}^{b}f(x)\,\mathrm{d}x=(b-a)f\left(\frac{a+b}{2}\right)+\frac{(b-a)^{3}}{24}f''(\xi).
\]

\item 设
\[
F(x)=\begin{cases}\displaystyle\int_{0}^{x}\cos\frac{1}{t}\,\mathrm{d}t, & x\neq 0,\\[2mm] 0, & x=0.\end{cases}
\]
求 $F'(0)$.

\item 函数 $u$ 在 $[0,+\infty)$ 上满足积分--微分方程
\[
\begin{cases}
u'(t)=u(t)+\displaystyle\int_{0}^{1}u(x)\,\mathrm{d}x,\\
u(0)=1.
\end{cases}
\]
试确定函数 $u$.
\end{exblock}
